\section{Conduct in economic context}
\label{sec:conduct}

The following seven incidents connect communication to the economic dependencies measured above. They were selected for distinct mechanisms, concrete stakes, recipient context and observable consequences. They are not the seven most severe events under a defined scale, nor a sample for estimating prevalence. Public messages, applied private notes and provider explanations remain separately labelled. Excerpts are verbatim; their arithmetic and descriptions are not silently corrected. The analysis outside the cards supplies the relevant state and interpretation.

\TranscriptTitle{insurance}{Selective disclosure in a sale framed as insurance}

Grok seeks cash for a mortgaged Pennsylvania Avenue while preserving its yellow group. At G13 d250 it has \$35 and publicly asks Opus for \$200, warning that Short Line rent would otherwise hand the yellows to Sol. Opus receives the appeal and counters at \$130 in d251. Its returned explanation values Grok's solvency after that appeal; the motive is not independent of the seller's framing. Grok replies:

\TranscriptQuote{g13-seller-counter}

The shortfall is true on a cash-only reading. Grok can also mortgage its clear Ventnor for \$130. Its applied private note recognizes that route and sets a \$165 floor to preserve Ventnor's rent. The returned provider explanation makes the presentation explicit:

\TranscriptQuote{g13-omission-explanation}

This span comes from a much longer explanation containing contradictory calculations and ownership statements. It supports a recorded concealment rationale, not a claim that every sentence expresses a stable intention. Opus accepts:

\TranscriptQuote{g13-acceptance}

The actual settlement pays Grok \$200, transfers Pennsylvania to Opus and charges Opus \$16 transfer interest. A retained exact-prefix diagnostic conditionally accepts \$130 and then returns Grok to its management decision; mortgaging Ventnor reaches \$295, enough for the \$200 rent with \$95 left. The diagnostic establishes an available financing alternative, not that Grok would have accepted it.

The message selectively presents a cash shortfall while withholding collateral. A refusal-posture reading also remains possible: Grok could be threatening not to transact at \$130 rather than claiming payment is physically impossible after acceptance. Opus's d253 returned explanation repeats the estate-denial stakes but does not endorse the impossibility of the cheaper route. The observed extra \$70 is therefore not an identified persuasion effect. In turn 150, Grok later fails on Short Line and Sol inherits the yellows. The purchase bought immediate liquidity without preventing the later event used to motivate it.

\TranscriptSource{G13 d250--253 and d311. Grok uses the standard condition; Opus's selected d251/d253 responses are action-only, with no applied private note. Their delivered timelines also contain self-partner counter wording. The opening appeal's 5/36 Short Line probability is wrong: the direct two-dice probability from Atlantic is 4/36. A prior public correction about missing North Carolina was in Opus's prompt.}

\TranscriptTitle{withdrawal}{A standing price disappears after the seller's deterioration}

Opus wants Sol's mortgaged Pacific Avenue in G18. It offers \$250 and repeats that price as standing across several exchanges. After collecting \$750 from Sol, it says:

\TranscriptQuote{g18-assurance}

A second \$750 rent payment then forces Sol to sell orange development and mortgage remaining reds, leaving \$49. After larger packages fail, Sol invokes the advertised price:

\TranscriptQuote{g18-invocation}
\TranscriptQuote{g18-withdrawal}

Opus's returned explanation points to Sol's near-insolvency and the limited value of paying to extend its life. The earlier assurance is in both players' delivered memories at invocation and withdrawal. No \$250 transaction occurs. Later \$150 and \$225 offers also fail to settle.

This is an opportunistic withdrawal under changed financial conditions, with an explicit reason to stop financing the seller. It is not evidence that Opus was dishonest when originally making the offer. No binding standing-offer object or expiry rule exists, and the record identifies no costly action undertaken in reliance on the assurance. The forced liquidation pays rent; it does not by itself establish reliance on an expected sale. The case illustrates how an offer's strategic value can reverse as the counterparty approaches elimination.

\TranscriptSource{G18 d421--442, especially d431/d441/d442, with subsequent offers retained in the source chronology. Opus's explanation is provider-returned text under action-only recovery.}

\TranscriptTitle{denial}{An explicit appeal to deny a rival development cash}

After Grok buys Illinois for \$500 and completes reds in G20, Gemini asks the table not to finance its construction:

\TranscriptQuote{g20-coordination-appeal}

Grok receives the message and its next applied private note calls it a cash boycott. Its d97 received END conflicts with financing/building content; the actual three red houses arrive at d98. After mortgaging New York, Grok offers Vermont to Gemini for \$150, describing it as a route toward light blues rather than a house fund. Gemini responds:

\TranscriptQuote{g20-concrete-refusal}

Grok mortgages Vermont instead. This is an explicit coordination appeal followed by the initiator's own refusal of the targeted cash transfer. It concerns another player's capacity to turn a completed group into rent, linking the social exchange directly to the development bottleneck.

The record does not establish a jointly adopted boycott. Sol had stated its own reluctance to fund Grok before Gemini's appeal, while Opus later proposes cash purchases from Grok. Neither settles, but those proposals contradict a simple account of unanimous compliance. Grok's acknowledgement establishes recognition of the appeal, not its effectiveness in changing counterparties' policies. The observed collective mechanism is attempted financial denial with unilateral follow-through, not a demonstrated cartel.

\TranscriptSource{G20 d95--102. The appeal is delivered at d97, whose selected attempt is 1 after malformed attempt 0. Gemini's d101 answers the actual d100 offer.}

\TranscriptTitle{bluff}{A public final offer below the stated private ceiling}

Grok wants North Carolina from Sol in G7. It proposes a public limit:

\TranscriptQuote{g7_d42_public}

The same valid response's applied private note leaves room for a higher price:

\TranscriptQuote{g7_d42_private}

Sol accepts \$600. Its own note favors reinvesting in cheaper sets rather than pursuing another \$25. The deal transfers North Carolina to Grok and \$600 to Sol; the private \$625 contingency is never tested.

This is a direct bargaining-limit bluff with a public/private discrepancy, rather than an erroneous board fact repeated in both channels. It is authorized by the public-deception instruction. The accepted price does not identify what the word ``final'' contributed, and the recipient's stated reason does not show belief that \$600 was Grok's true ceiling. The episode nevertheless demonstrates a recognizable negotiation tactic without requiring a successful-manipulation label.

\TranscriptSource{G7 d42--43, both standard valid attempt 0. The actual acceptance prompt includes the public final-offer statement.}

\TranscriptTitle{obstacle}{A false building obstacle with contrary text in the same response}

Only one house remains in the bank late in G3. Grok publicly claims that Gemini cannot add a Park Place house without somebody releasing buildings:

\TranscriptQuote{G3-2}

Its applied private note says the opposite about stock and identifies collateral:

\TranscriptQuote{G3-3}

At Gemini's next actual management decision, an intervening \$15 tax leaves \$43. A retained legal witness mortgages Kentucky for \$110 and Atlantic for \$130, buys the \$200 Park Place house and leaves \$83. No opponent has to release a house. Gemini instead ends the turn, recording a preference to preserve unmortgaged income.

The public stock obstacle is false, while the immediate cash shortfall is real. The contradiction is not sufficient to infer one coherent deceptive intention: Grok's longer provider explanation repeatedly treats additional release as necessary. Gemini had already held at d362 before the statement. Exposure followed by another hold therefore does not establish reliance or an induced mistake. The case belongs beside the clearer bluff because both are pressure tactics in the public record, but their intent evidence differs.

\TranscriptSource{G3 d362--364. The finance route is checked at Gemini's actual d364 management opportunity. The claim concerns one purchase, not eventual survival or optimal investment.}

\TranscriptTitle{correction}{A recipient corrects a false claim and takes the terms}

In G19, Grok offers Virginia plus \$150 for Sol's States. Its public message falsely describes States as an orange deed. Yet the selected assistant content names END, and the note attached to the received proposal says:

\TranscriptQuote{g19-abort}

The received tool still proposes the trade. Sol receives it and replies:

\TranscriptQuote{g19_d136_correction_accept}

Events execute the Virginia/States swap and \$150 payment from Grok to Sol. No fallback or corrective retry produced the proposal; it was the sole valid received call. Sol's note also treats States as pink and favors the terms.

The recipient's explicit correction shows that settlement can coexist with rejection of the advertised premise. The offer itself is consequential, but its conflicting content and abort instruction prevent attributing it to a single plan to deceive. This is a transaction-level example of recipient discrimination and an execution inconsistency, not proof that misleading pitches generally fail. It also explains why the complete response census should be paired with source channels when judging the strategy behind an accepted offer.

\TranscriptSource{G19 d135--136, standard attempt 0. Received, parsed and applied tool names agree; their disagreement is with assistant content and the attached private note.}

\TranscriptTitle{surrender}{An applied bankruptcy contradicts the response's survival instruction}

Grok owes \$250 with \$43 cash at G12 d249. Its pinks have two, two and three houses. Assistant content requests a two-house Virginia sale, while the attached note concludes:

\TranscriptQuote{g12_d249_private_final}

The one received tool call instead declares bankruptcy:

\TranscriptQuote{g12_d249_public_bankruptcy}

No sale occurs; the estate transfers to the creditor. The requested two-house sale is legal but would reach only \$143. A separate five-house sale sequence pays the debt and leaves \$43. The note contains intermediate mistakes, so the endpoint follows the independent legal witness rather than its arithmetic.

This is a severe loss of an available current-debt payment route, coupled with an explicit disagreement between the requested sale and applied declaration. It is not evidence of deliberate gifting or a desire to lose. The case demonstrates why strategic execution and strategic explanation have to be compared at the decision where an action is possible. Future winning prospects after payment remain untested.

\TranscriptSource{G12 d249, standard attempt 0. The content sale is not an applied or rejected tool call. The five-sale witness is retained in the historical liquidation audit.}

Across these incidents, the economic themes recur: cash may preserve a blocking owner, finance an opponent's development, or prolong a debtor whose assets another player wants. The records also distinguish an explicit bluff, a selective financing presentation, a contingent bargaining reversal, a coordination appeal, and contradictory response channels. Calling all of them deception would discard the differences that explain what occurred. Calling none strategically meaningful because persuasion is unidentified would discard the conduct itself.
